\section{Proposed Method: Source Facts and Semantic Feedback}
\label{sec:method}
\subsection{Design hypothesis and implementation boundary}
The proposed method strengthens the same LLM by supplying evidence that it can
use to revise a candidate and its mapping. It combines a source analysis stage,
an obligation ledger, claim checking, and bounded revision. This is an integration
proposal, not a claim that the full architecture is already implemented. The
existing infrastructure supports source coordinates, structured predictions,
private review records, and trusted local checks. A separate prototype accepts
structured vertex-cover QUBO templates and returns development-test feedback.
Automatic extraction of full program contracts, general numeric-state analysis,
and production-program transformation remain outside that implemented prototype.

The technical hypothesis is specific: errors in a mapping often arise because
a mathematical pattern is detached from its source preconditions or observable
selection rules. A dependency fact or a reachable counterexample can correct that
detachment. This hypothesis is narrower than assuming that more tools or a longer
prompt improve the entire task. The ablations in Section~\ref{sec:evaluation}
are designed to isolate which evidence, if any, helps.

\subsection{Candidate-centered program facts}
Given a nominated region, the proposed analysis constructs a bounded dependency
slice covering live inputs, state updates, guards, and uses of its result. A fact
has a source location, an analysis origin, an assumption set, and a confidence
status. Examples include a loop's iteration order, an explicit finite-value check,
the use of a tuple in an ordered report, and a caller-selected iteration budget.
The model sees these facts with their limitations; the analyzer must not label a
heuristic inference as a theorem.

For dynamic Python features outside the supported analysis, the method returns
an unresolved dependency rather than assuming purity or a closed call graph.
Alias-sensitive updates and exception paths require conservative treatment.
Facts established at the program entry are not silently lifted to every later
program point. In particular, finite numeric inputs do not imply that intermediate
floating point operations remain finite. A dynamic witness can demonstrate a
reachable state; a collection of successful traces cannot prove that an unobserved
state is unreachable.

An obligation ledger links these facts to the proposal. For each obligation it
stores the required behavior, source anchors, proposed implementation or proof,
available check, and status. The ledger is not a hidden answer key: public facts
may be derived from the input, whereas evaluator-only reference judgments remain
withheld. The full method is compared with both an identical prompt-only ledger
and a fact-assisted workflow to avoid crediting the presentation format alone.

\subsection{Checking mapping claims}
The model first supplies the original conditional plan. Supported parts may then
be expressed in a restricted claim representation: finite predicates, coefficient
templates, variable mappings, optimization direction, decoding, and preconditions.
For a plan originally written in prose, transcription is logged and checked against
the source passage. If a reviewer must invent missing coefficients or assumptions,
that addition is a separate proposal and not credited to the model.

The vertex-cover prototype illustrates the representation boundary. It interprets
a whitelist of arithmetic expressions, builds a quadratic polynomial, and uses
exact rational arithmetic. Templates describe a family of graphs; they cannot
select answers by looking up final-test instances. Variable permutation, bit
complementation, constant offsets, and appropriate scaling may yield equivalent
solutions and must be accepted. A formula outside the representation's expressivity
receives a representation outcome, not a universal judgment that no valid mapping
exists.

For each supported claim $q$, the checker returns
\begin{equation}
 V(q)=(s,\mathcal{D},\alpha,w),\quad
 s\in\{\text{supported},\text{refuted},\text{unknown}\},
 \label{eq:verdict}
\end{equation}
where $\mathcal{D}$ is the checked domain, $\alpha$ the assumptions, and $w$ an
evidence record. A refutation includes a witness and the violated obligation;
finite support includes the actual test boundary. A universally quantified
correctness claim requires a proof beyond a finite successful test run.

\subsection{Bounded feedback to the same model}
Feedback exposes the first failing development instance, the claimed output or
decoded optimum, the required observation, and the violated obligation. It does
not expose a replacement reference formula or any final-test result. When all
development checks pass, feedback says only that the declared checks passed.
The same model revises its whole plan within a fixed call budget, retaining the
ability to narrow a claim, state an unresolved obligation, or recommend keeping
the original classical routine.

Table~\ref{tab:procedure} describes the proposed orchestration. The feedback loop
resembles the candidate/verification separation of counterexample-guided
synthesis~\cite{sketching}, but an LLM proposal plus bounded testing does not
inherit synthesis completeness or termination guarantees. The workflow ends on
the predetermined budget, not only when a favorable answer appears.

\begin{table}[tb]
\caption{Proposed bounded workflow. Evaluator-reserved checks never produce
feedback within a measured run.}
\label{tab:procedure}
\begin{tabularx}{\linewidth}{lY}
\toprule
Step & Operation and retained evidence \\
\midrule
Input & Freeze public program, contract, model configuration, and budgets. \\
Propose & Save raw candidate regions, eligibility, mapping, recommendation, and open obligations. \\
Analyze & Derive source facts with assumptions and source anchors; record unsupported constructs. \\
Check & Validate supported claims against development checks; retain scope and witnesses. \\
Revise & Send allowed facts and feedback to the same model; save every revision. \\
Stop & End at the fixed budget or protocol stopping condition, retaining failures. \\
Evaluate & Run reserved checks and blinded review once on the designated final artifact. \\
\bottomrule
\end{tabularx}
\end{table}

\subsection{Reachable states and classical fallback}
A useful diagnostic concerns a proposed first-maximum predicate. Over ordinary
totally ordered finite values, ``no smaller than every candidate, with no earlier
equal value'' can characterize a first maximum. With NaN values, these comparisons
do not reproduce a sequential incumbent-update rule. A counterexample is relevant
to a program only if the state is permitted or reachable from a permitted input.
The planned numeric analysis therefore binds a witness to the original program
point and execution trace instead of testing an arbitrary impossible array.

This distinction also constrains feedback scope. Refuting a predicate does not
refute every conditional plan that mentions it. A plan may detect an exceptional
state and fall back to the classical routine. Review must determine whether that
guard is actually supplied, occurs before the incorrect operation, and preserves
the full contract. A vague promise to ``verify and fall back'' is incomplete until
the verification condition and fallback behavior are specified. Conversely, an
exact but expensive fallback may preserve correctness while eliminating any
practical benefit. The method records these dimensions separately.

\subsection{Resources and adoption}
Structural mapping and practical adoption use different evidence. A resource
ledger records state preparation, reversible oracle construction, ancilla cleanup,
logical and decomposed circuit resources, repeated measurements, classical
optimization, data transfer, checking, and fallback. Unknown terms remain unknown.
Any circuit counts name the representation and compilation target; they are not
used as an unexplained profitability score.

The method never turns simulator wall time into a quantum speedup claim. For this
paper, its own usefulness is measured as improved analysis or mapping quality and
the classical/model cost of obtaining it. An end-to-end advantage study would
need implemented hybrid programs, appropriate classical competitors, hardware
assumptions, and a separate experimental protocol.
